\section{Signed lattice symmetry of the quadratic precision}
\label{sec:quadratic-lattice-symmetry}
The path attachment uses the quadratic precision of the selected bed.
We verify its invariance under signed permutations of the coordinate axes,
including the change of anchor caused by reversing a bond. This is a
property of the bilinear form on one-cochains; the curvature metric need
not satisfy a stronger symmetry away from exact curvatures.

Let $s$ be a signed permutation of the four axes, acting on sites and
momenta by the same orthogonal integer matrix. A positive bond $(x,a)$
becomes the directed bond from $sx$ to $sx+s e_a$. If $s e_a=-e_b$,
its value in positive-bond coordinates is minus the value at $(sx-e_b,b)$.
This defines the pullback $R_s$ on real one-cochains. For a canonical
plaquette $p=(a,b)$, $a<b$, let $r_s(x,p)$ be its image with its two
axes reordered increasingly and its lower vertex shifted for each
reversed axis. Its orientation sign is
\[
 \epsilon_s(p)=\epsilon_s(a)\epsilon_s(b)
                   \operatorname{sgn}(s_{\mathrm{perm}}(b)-s_{\mathrm{perm}}(a)),
\]
where $\epsilon_s(a)\in\{-1,1\}$ is the sign of the image axis.
In the permutation subscripts, only the unsigned permutation is meant.
Equivalently the last factor is $1$ if the two image axes are in
increasing order and $-1$ otherwise. The map $r_s$ is a permutation of
canonical plaquettes on every finite torus and $\epsilon_s(p)^2=1$.

Let $d$ denote the discrete exterior derivative. The boundary of the
image of an oriented plaquette is the image of its boundary, with the
same orientation sign, so
\begin{equation}\label{eq:quadratic-curl-rotation}
       (dR_sa)(x,p)=\epsilon_s(p)(da)(r_s(x,p)).
\end{equation}
One can verify the identity directly by substituting
$(da)(x,a,b)=a(x,a)+a(x+e_a,b)-a(x+e_b,a)-a(x,b)$.
Reversing an axis replaces each affected bond by its negative at the
shifted lower endpoint; exchanging the two axes reverses the whole
expression. These two operations generate all the signs and permutations,
so the formula holds also across the torus seam and for even periods.

Assume the blocking penalty is positive and the kernel level is at
least four, as in the retained construction. Let $P$ be the unamended curvature row from
Section~\ref{sec:positive-row-producer}. Its Fourier symbol is diagonal
in the plaquette axes. For $q\in[-\pi,\pi]^4\setminus\{0\}$, its entry at
$p=(a,b)$ is
\[
 \widehat P_{pp}(q)=
      \rho(q)^{-1}D_a(q)^{-1}D_b(q)^{-1},
 \qquad
 \rho(q)=\sum_c\frac{4\sin^2(q_c/2)}{D_c(q)}.
\]
Here $D_c$ are the positive direction denominators of the selected
finite tower. Their defining alias sums satisfy
$D_{s_{\mathrm{perm}}(c)}(sq)=D_c(q)$: a signed permutation bijects the
alias indices, preserves their scalar product weights, and permutes
the single distinguished axis. The identity also follows term by term
through the finite tower recursion. The sine squares are unchanged
by sign and permuted with the same axes. Consequently $\rho(sq)=\rho(q)$
and
\[
            \widehat P_{pp}(q)
                =\widehat P_{p',p'}(sq),\qquad p'=s p
\]
with the image axes canonically ordered. At $q=0$ the unamended symbol is the identity matrix, so the same
identity holds there too. Off-diagonal entries vanish.

The Laurent rows are absolutely summable. Uniqueness of Fourier
coefficients and the change of variable $q\mapsto sq$ therefore give
\[
 P_{pp}(z)=P_{p'p'}(sz),\qquad P_{pq}=0\quad(p\ne q).
\]
Let $P_n$ denote the periodized row on a torus of side $n$.
The lower vertex of an image plaquette has the form $sx+t_p$, where
$t_p$ depends only on its axes. For two plaquettes with the same axes,
this shift cancels from their separation. Reindexing the absolutely
convergent periodization sum by $z\mapsto sz$ thus proves
\begin{equation}\label{eq:quadratic-face-symmetry}
 P_n((x,p),(y,q))\epsilon_s(p)\epsilon_s(q)
       =P_n(r_s(x,p),r_s(y,q)).
\end{equation}
For unequal axes both sides are zero; for equal axes the two signs
multiply to one. This cancellation is why the anchor shifts cannot
be ignored before passing to the bilinear form.

Using \eqref{eq:quadratic-curl-rotation} and
\eqref{eq:quadratic-face-symmetry}, expand the finite double sum:
\begin{align*}
 \langle dR_sa,P_n dR_sb\rangle
 &=\sum_{u,v}\epsilon_s(u)\epsilon_s(v)
          (da)(r_su)P_n(u,v)(db)(r_sv)\\
 &=\sum_{u,v}(da)(r_su)P_n(r_su,r_sv)(db)(r_sv)\\
 &=\langle da,P_n db\rangle.
\end{align*}
The last equality reindexes by the permutation $r_s$.

The positive amendment of Section~\ref{sec:curvature-horizon-producer}
vanishes on exact curvatures: if $M$ is the amended metric, then
$d^*M_nd=d^*P_nd$. Hence its quadratic bilinear form satisfies
\[
          \langle R_sa,d^*M_ndR_sb\rangle
                        =\langle a,d^*M_ndb\rangle.
\]
For a Lie-algebra-valued field, expand in an orthonormal colour basis.
The signed lattice pullback acts separately in every colour, and the
gauge quadratic form is the sum of the real bilinear forms just
considered. Summing their equalities proves the required invariance
for every pair of such fields and every period. The positive square
root used to represent the amended metric is not needed in this
symmetry argument.
